\Needspace{12\baselineskip}
\originalchapter{Lab4：分子动力学的全部公开任务}
原Lab4截止于2005年4月21日，主体为Ni-EAM积分器比较、Kr熔化、规模与温度涨落。此前约化LJ两步长、力差分和MC验证已经独立通过，本章引用该证据，不更改或把它们改名为Ni/Kr原软件运行。原GULP与MOLDY执行、材料参数、单位和积分器实现须分别记录。

\originalsection{三类积分器的同条件比较}
\begin{exercise}\label{task:lab4-1}
（原题1，五项）用GULP的Ni-EAM、300 K：(A)对leapfrog Verlet扫描时间步，画平均总能量；(B)同图加入velocity Verlet和Gear predictor--corrector；(C)画三者能量标准差随步长；(D)在.001 ps画瞬时能量；(E)评价稳定性、准确性和计算代价。
\end{exercise}
\begin{solution}
必须给三算法相同初始坐标、物理初始速度、势、胞、固定\emph{物理时长}及平衡/生产分段；不能用相同步数比较而改变实际观测时间。leapfrog速度位于半时刻，velocity Verlet在整时刻；计算动能时若不插值到同一时刻，就能产生平均能量偏差。这是读出约定问题，不能自动归结为轨迹不稳定。

原公开答案说明：Verlet族在较大步长前稳定；Gear在其小步长算例的漂移/涨落表现较好，但较大步长时更早失稳；leapfrog表观平均能量差与错开的速度时刻有关。它没有给一种算法对所有势、步长普遍最优的结论。Gear多阶预测和力校正的代价依具体实现，不能只按公式项数排行。

报告应分别给 $\bar E-E_{\rm reference}$、$\operatorname{std}E$、拟合漂移斜率、最大振动频率与CPU时间。能量标准差小不等于没有系统偏差，恒温器也可能掩盖积分漂移。已通过的\emph{约化LJ模型}两步长结果为 $.00031081,.00007766$，支持该轨迹下二阶误差尺度；它仍不是原Ni-EAM三算法图。原公开脚本确有Ni的Sutton--Chen/EAM系数，GULP尚未执行的部分保留完整方法，未虚构图点。
\end{solution}

\originalsection{Kr熔化与两相界面}
\begin{exercise}\label{task:lab4-2}
（原题2）(A)以LJ模型在多温度直接加热固态Kr，自选并论证时间步、胞大小与系综，保存平衡/采样时间；用均方位移等多种分析确定熔化，讨论其他易理解但需新分析代码的指标及加热法的高低偏差。(B)阅读原题Morris等所列论文，描述两相模拟法，并据此计算Kr熔化温度；本文不复制第三方论文全文。
\end{exercise}
\begin{solution}
加热固态时保存展开坐标以求MSD，比较弹道、振动平台和真正长时线性区；检查径向分布、结构因子、键角/局部晶序、能量与体积突变及速度相关。单看一个快照或短时MSD不能可靠定温。NVT、NPT、NVE改变所约束的变量，尤其压力和体积松弛会影响熔化；没有明确势参数与状态点时不能给“Kr唯一熔点”。

原官方答案的例子：试过.001--1 ps，超过.01 ps出现不稳定碰撞；直接加热快照估计130--150 K，MSD和势能突变定位在140--150 K。更大 $7^3$ 超胞、5000平衡步和45000生产步的NPE例子给约110--120 K，并指出其系综实现限制了液气转变。以上是不同条件的原结果，不能挑一个数代表同一协议，也不能把约化LJ的温度.3直接改名为这些Kelvin结果。

两相法准备一半固体、一半液体及足够分离的界面，在同一P、组成和受控边界下观察界面方向；高于熔点液体增长、低于熔点固体增长。NVE或适当定压定焓协议中，潜热释放/吸收还反馈温度，不能把初始温度等同最终温度。原答案的NVE例子给初始100、130、140 K分别到119.5、120.9、138.4 K；前两例接近约120 K，第三例仍需观察界面是否停止。有限体系界面移动、体积、时长与局部晶序都要保存，才能确定误差区间。

单相加热通常过热、冷却可过冷，因为形核势垒需要涨落触发；两相法预置界面减少形核等待。原答案中“总会高估/低估”是该常见协议的趋势，边界、缺陷和扫描速率可改变它，不应作为普遍严格定理。本项目随后已真实运行原Kr势与晶体输入；两相输入的直接启动失败及进一步准备后的有限轨迹见后文“原材料输入的实际计算记录与解释”。它们尚不能独立给平衡熔点，原官方结果与本次轨迹分开保留。
\end{solution}

\originalsection{规模与温度涨落}
\begin{exercise}\label{task:lab4-3}
（原题3）周期FCC超胞：(A) $n^3$ 常规胞含多少原子；(B)全部配对和截断各需多少力/距离计算；(C)按CPU成本比较HF、LJ、EAM和DFT并解释；(D)如何减少距离计算？
\end{exercise}
\begin{solution}
已核验片段的完整推导保留：$N=4n^3$，无序对 $N(N-1)/2$；固定密度截断内平均z个邻居给 $Nz/2$ 个相互作用，但若仍先计算所有距离，候选筛选仍是二次规模。胞链表把候选范围局部化，带皮层邻居表允许多步复用；累计位移、胞形变化和截断需一致检查。平方距离先筛选后求根，成对反力只算一次，不改变物理相互作用。

原预期一般排行为LJ、EAM、DFT、HF从较低到较高成本；基组、交换处理、电子数、稀疏性及硬件会改变它。算法加速仍须和未加速小体系核对能量与力，不能以速度快作为准确性的证据。
\end{solution}

\begin{exercise}\label{task:lab4-4}
（原题4）(A)取题1 velocity Verlet、.001 ps的生产期瞬时温度曲线；(B)用常规胞边数4、5、6等增大体系、同算法和步长，在同图比较；(C)解释涨落变化；(D)导出温度标准差与原子数的关系，并用(B)验证。
\end{exercise}
\begin{solution}
移除整体动量后自由度 $f=3N-3$，瞬时动温 $T_{\rm kin}=2K/(f k_B)$。在正则分布、独立二次自由度下，$2K/(k_BT)$ 服从 $\chi_f^2$，所以 $\operatorname{Var}(T_{\rm kin})=2T^2/f$，标准差相对均值为 $\sqrt{2/f}\propto N^{-1/2}$。这给理论比较基准，不等于本原题生产期一定处于正则系综。

微正则总能量固定使K与势能反相关，动力温度涨落还含热容约束；不能把NVE数据的系数强制拟合成正则公式。设热容 $C_V$ 为全体系值，常见微正则关系在适用的热力学近似下为 $\operatorname{Var}K/\langle K\rangle^2\simeq(2/f)[1-f k_B/(2C_V)]$。若体系接近相变、存在慢模式或整体温漂，这些前提须检验。

原任务的验证应按相同生产时长或足够独立采样、相同系综和初温，画相对标准差对 $1/\sqrt{3N-3}$，给块误差和相关时间；不可把高频采样数量当作独立样本数。原Ni-EAM温度轨迹未生成，完整原任务与理论、文件列和拟合方法保留，不以解析方差假装实跑数据。
\end{solution}

\originalsection{全部公开加分路径}
\begin{exercise}\label{task:lab4-extra-md}
（原加分题）从零写MD：基础含力、积分器、位置输出；进一步含构型输入输出、K/U/E、周期边界、速度重标定平衡；再含邻居表、压力/温度、MSD、径向分布；最高阶包括其他系综、胞链表、多势和MPI并行。原题按这四层各5分。
\end{exercise}
\begin{solution}
本卷已运行代码覆盖成对力、velocity Verlet、K/U/E、周期最小镜像、温度自由度及展开坐标的基础实现，力差分和两步长已独立核验；并未宣称邻居表、MPI或可变胞系综均已完成。复用时将势接口、积分状态和观测量输出分开，加入功能后先在小体系核对能量/力，再检验采样系综。题目的高级层级作为完整要求保留，不能因本模型暂不需要而删掉。
\end{solution}
\begin{exercise}\label{task:lab4-extra-moldy}
（原替代编程题）修改MOLDY，使其除成对势外支持pair-functional/EAM类势。
\end{exercise}
\begin{solution}
需先计算各原子的密度 $\rho_i=\sum_jf(r_{ij})$ 和嵌入导数，再对每对构造 $\phi'(r)+[F'(\rho_i)+F'(\rho_j)]f'(r)$ 的力，保持Newton反力及能量计数。单程仅替换v(r)不够，因为一处位移改变多个环境密度。测试应含非对称位移的能量差分、截断平滑与完整跨边界邻居。本项目未修改原MOLDY源码，实施层面的结果保持未执行。
\end{solution}
\begin{exercise}\label{task:lab4-extra-ni}
（原另一加分题）用GULP MD对给定Ni-EAM直接加热，研究熔化。
\end{exercise}
\begin{solution}
用题2的多观测量和加热形核偏差分析，但材料、EAM参数、时间单位、压力与晶格均改为Ni原输入；不能将Kr或本卷约化LJ的温度数字平移为Ni结果。保存每温度段的平衡/生产和重复初态；原Ni输入已实际取得、GULP未执行的缺口如实列出，题面和可复用流程完整保留。
\end{solution}
